\documentclass[11pt,a4paper]{article}

\usepackage[utf8]{inputenc}
\usepackage[T1]{fontenc}
\usepackage[spanish]{babel}
\usepackage{amsmath,amssymb,mathtools}
\usepackage{geometry}
\usepackage{lmodern}
\usepackage{microtype}

\geometry{margin=2.4cm}
\setlength{\parindent}{0pt}
\setlength{\parskip}{0.45em}

\newcommand{\R}{\mathbb{R}}
\newcommand{\sen}{\operatorname{sen}}
\newcommand{\Dom}{\operatorname{Dom}}

\begin{document}

\begin{center}
{\LARGE\bfseries Análisis Matemático II}\\[0.4em]
{\Large Tema 2: Límites y continuidad en funciones de varias variables}\\[0.4em]
{\large Soluciones}
\end{center}

\section{Ejercicio 1}
\subsection{Solución}

Queremos demostrar:

\[
\lim_{(x,y)\to(0,0)}
\frac{x^2y^2}{x^2+y^2}
=
0.
\]

Se tiene:

\[
x^2y^2
\leq
\frac{(x^2+y^2)^2}{4}.
\]

Por tanto:

\[
0
\leq
\frac{x^2y^2}{x^2+y^2}
\leq
\frac{x^2+y^2}{4}.
\]

Si:

\[
0<\sqrt{x^2+y^2}<\delta,
\]

entonces:

\[
\left|
\frac{x^2y^2}{x^2+y^2}
\right|
<
\frac{\delta^2}{4}.
\]

Tomando:

\[
\delta=2\sqrt{\varepsilon},
\]

se obtiene:

\[
\left|
\frac{x^2y^2}{x^2+y^2}
\right|
<
\varepsilon.
\]

Por tanto:

\[
\boxed{
\lim_{(x,y)\to(0,0)}
\frac{x^2y^2}{x^2+y^2}
=
0
}
\]

\section{Ejercicio 2}
\subsection{Solución}

Sea:

\[
f(x,y)
=
\frac{xy^3}{x^2+y^6}.
\]

Según una recta:

\[
y=mx,
\]

se obtiene:

\[
f(x,mx)
=
\frac{m^3x^4}{x^2+m^6x^6}
=
\frac{m^3x^2}{1+m^6x^4}.
\]

Así:

\[
\boxed{
\lim_{x\to0}f(x,mx)=0
}
\]

para todo \(m\in\R\).

Según el subconjunto:

\[
x=y^3,
\]

se tiene:

\[
f(y^3,y)
=
\frac{y^6}{y^6+y^6}
=
\frac12.
\]

Por tanto:

\[
\boxed{
\lim_{y\to0}f(y^3,y)=\frac12
}
\]

Como se obtienen límites distintos:

\[
\boxed{
\lim_{(x,y)\to(0,0)}
\frac{xy^3}{x^2+y^6}
\text{ no existe}
}
\]

\section{Ejercicio 3}
\subsection{Solución}

Consideramos:

\[
f(x,y)
=
\frac{(x+1)x+(2y-1)y}{x-y}.
\]

Según las rectas:

\[
y=mx,
\]

con \(m\neq1\):

\[
f(x,mx)
=
\frac{x^2+x+2m^2x^2-mx}{x-mx}
\]

\[
=
\frac{(1-m)+(1+2m^2)x}{1-m}.
\]

Por tanto:

\[
\boxed{
\lim_{x\to0}f(x,mx)=1
}
\]

para toda recta \(y=mx\), \(m\neq1\).

Ahora tomamos:

\[
y=x-\alpha x^2,
\qquad
\alpha\neq0.
\]

Entonces:

\[
x-y=\alpha x^2
\]

y:

\[
(x+1)x+(2y-1)y
=
(3+\alpha)x^2
-
4\alpha x^3
+
2\alpha^2x^4.
\]

Así:

\[
f(x,x-\alpha x^2)
=
\frac{3+\alpha}{\alpha}
-
4x
+
2\alpha x^2.
\]

Luego:

\[
\boxed{
\lim_{x\to0}
f(x,x-\alpha x^2)
=
1+\frac3\alpha
}
\]

que depende de \(\alpha\).

Por tanto:

\[
\boxed{
\lim_{(x,y)\to(0,0)}f(x,y)
\text{ no existe}
}
\]

\section{Ejercicio 4}
\subsection{Solución}

Sea:

\[
f(x,y)
=
\frac{xy-x+y}{x+y}.
\]

Primero:

\[
\lim_{y\to0}f(x,y)
=
\frac{-x}{x}
=
-1.
\]

Por tanto:

\[
\lim_{x\to0}
\left(
\lim_{y\to0}f(x,y)
\right)
=
-1.
\]

En el otro orden:

\[
\lim_{x\to0}f(x,y)
=
\frac{y}{y}
=
1.
\]

Así:

\[
\lim_{y\to0}
\left(
\lim_{x\to0}f(x,y)
\right)
=
1.
\]

Como los límites reiterados son distintos:

\[
\boxed{
\lim_{(x,y)\to(0,0)}
\frac{xy-x+y}{x+y}
\text{ no existe}
}
\]

\section{Ejercicio 5}
\subsection{Solución}

Realizamos el cambio:

\[
x=r\cos\theta,
\qquad
y=r\sen\theta.
\]

Entonces:

\[
\frac{x^2y}{x^2+xy+2y^2}
=
r
\frac{\cos^2\theta\,\sen\theta}
{\cos^2\theta+\cos\theta\sen\theta+2\sen^2\theta}.
\]

El denominador angular es estrictamente positivo:

\[
\cos^2\theta+\cos\theta\sen\theta+2\sen^2\theta
=
\left(
\cos\theta+\frac12\sen\theta
\right)^2
+
\frac74\sen^2\theta.
\]

Por tanto, el cociente angular está acotado y:

\[
\boxed{
\lim_{(x,y)\to(0,0)}
\frac{x^2y}{x^2+xy+2y^2}
=
0
}
\]

\section{Ejercicio 6}
\subsection{Solución}

Con:

\[
x=r\cos\theta,
\qquad
y=r\sen\theta,
\]

se obtiene:

\[
(x+y)\ln(x^2+y^2)
=
r(\cos\theta+\sen\theta)\ln(r^2).
\]

Así:

\[
=
2r\ln r\,(\cos\theta+\sen\theta).
\]

Como:

\[
|\cos\theta+\sen\theta|
\leq
\sqrt2
\]

y:

\[
\lim_{r\to0^+}r\ln r=0,
\]

se concluye:

\[
\boxed{
\lim_{(x,y)\to(0,0)}
(x+y)\ln(x^2+y^2)
=
0
}
\]

\section{Ejercicio 7}
\subsection{Solución}

Sea:

\[
f(x,y)
=
\frac{(1-\cos(xy))\sen x}{x^2+y^2}.
\]

Utilizamos:

\[
1-\cos(xy)
\sim
\frac{x^2y^2}{2}
\]

y:

\[
\sen x\sim x.
\]

Por tanto:

\[
f(x,y)
\sim
\frac{x^3y^2}{2(x^2+y^2)}.
\]

Además:

\[
\left|
\frac{x^3y^2}{2(x^2+y^2)}
\right|
\leq
\frac{(x^2+y^2)^{5/2}}
{2(x^2+y^2)}
=
\frac12(x^2+y^2)^{3/2}.
\]

Luego:

\[
\boxed{
\lim_{(x,y)\to(0,0)}
\frac{(1-\cos(xy))\sen x}{x^2+y^2}
=
0
}
\]

\section{Ejercicio 8}
\subsection{Solución}

Sea:

\[
f(x,y)
=
\frac{e^{xy}-1}
{\sen x\,\ln(1+y)}.
\]

Utilizamos los infinitésimos equivalentes:

\[
e^{xy}-1\sim xy,
\]

\[
\sen x\sim x,
\]

\[
\ln(1+y)\sim y.
\]

Por tanto:

\[
f(x,y)
\sim
\frac{xy}{xy}
=
1.
\]

Así:

\[
\boxed{
\lim_{(x,y)\to(0,0)}
\frac{e^{xy}-1}
{\sen x\,\ln(1+y)}
=
1
}
\]

\section{Ejercicio 9}
\subsection{Solución}

Sea:

\[
f(x,y)
=
\frac{x^2+xy+y^2}{x^2+y^2}.
\]

Los límites reiterados son:

\[
\lim_{x\to0}
\left(
\lim_{y\to0}f(x,y)
\right)
=
1
\]

y:

\[
\lim_{y\to0}
\left(
\lim_{x\to0}f(x,y)
\right)
=
1.
\]

Según:

\[
y=mx,
\]

se obtiene:

\[
f(x,mx)
=
\frac{1+m+m^2}{1+m^2}.
\]

Por ejemplo:

\[
m=0
\Longrightarrow
1,
\]

\[
m=1
\Longrightarrow
\frac32.
\]

Por tanto:

\[
\boxed{
\lim_{(x,y)\to(0,0)}f(x,y)
\text{ no existe}
}
\]

En coordenadas polares:

\[
f(r,\theta)
=
1+\cos\theta\sen\theta,
\]

que depende de \(\theta\), confirmando el resultado.

\section{Ejercicio 10}
\subsection{Solución}

Con:

\[
x=r\cos\theta,
\qquad
y=r\sen\theta,
\]

se obtiene:

\[
\frac{x^3+y^3}{x^2-y^2}
=
r
\frac{\cos^3\theta+\sen^3\theta}
{\cos^2\theta-\sen^2\theta}.
\]

Factorizando:

\[
=
r
\frac{\cos^2\theta-\cos\theta\sen\theta+\sen^2\theta}
{\cos\theta-\sen\theta}.
\]

Para un \(\theta\) fijo distinto de \(\pi/4\), el límite es:

\[
0.
\]

Sin embargo, tomamos:

\[
\theta=\frac{\pi}{4}-r.
\]

Entonces:

\[
\cos\theta-\sen\theta
\sim
\sqrt2\,r
\]

y:

\[
\cos^2\theta-\cos\theta\sen\theta+\sen^2\theta
\to
\frac12.
\]

Por tanto:

\[
\lim_{r\to0^+}
r
\frac{\cos^2\theta-\cos\theta\sen\theta+\sen^2\theta}
{\cos\theta-\sen\theta}
=
\frac1{2\sqrt2}.
\]

Como se obtienen valores distintos:

\[
\boxed{
\lim_{(x,y)\to(0,0)}
\frac{x^3+y^3}{x^2-y^2}
\text{ no existe}
}
\]

\section{Ejercicio 11}
\subsection{Solución}

Sea:

\[
f(x,y)
=
\begin{cases}
\dfrac{x^2+2y^2}{2x^2+y^2},
&
(x,y)\neq(0,0),\\[2mm]
0,
&
(x,y)=(0,0).
\end{cases}
\]

Primero:

\[
\lim_{y\to0}f(x,y)
=
\frac{x^2}{2x^2}
=
\frac12.
\]

Así:

\[
\boxed{
\lim_{x\to0}
\left(
\lim_{y\to0}f(x,y)
\right)
=
\frac12
}
\]

En el otro orden:

\[
\lim_{x\to0}f(x,y)
=
\frac{2y^2}{y^2}
=
2.
\]

Por tanto:

\[
\boxed{
\lim_{y\to0}
\left(
\lim_{x\to0}f(x,y)
\right)
=
2
}
\]

Como son distintos:

\[
\boxed{
\lim_{(x,y)\to(0,0)}f(x,y)
\text{ no existe}
}
\]

\section{Ejercicio 12}
\subsection{Solución}

Sea:

\[
f(x,y)
=
\begin{cases}
y\sen\left(\dfrac1x\right)
+
x\sen\left(\dfrac1y\right),
&
x\neq0,\ y\neq0,\\[2mm]
0,
&
x=0\text{ o }y=0.
\end{cases}
\]

Para \(x\neq0\):

\[
\lim_{y\to0}
x\sen\left(\frac1y\right)
\]

no existe.

Por tanto:

\[
\boxed{
\lim_{x\to0}
\left(
\lim_{y\to0}f(x,y)
\right)
\text{ no existe}
}
\]

Análogamente:

\[
\boxed{
\lim_{y\to0}
\left(
\lim_{x\to0}f(x,y)
\right)
\text{ no existe}
}
\]

Esto no permite concluir que el límite doble no exista.

De hecho:

\[
|f(x,y)|
\leq
|x|+|y|.
\]

Como:

\[
|x|+|y|
\to0,
\]

se obtiene:

\[
\boxed{
\lim_{(x,y)\to(0,0)}f(x,y)=0
}
\]

\section{Ejercicio 13}
\subsection{Solución}

Sea:

\[
f(x,y)
=
\frac{x\sen y}{y\sen x}.
\]

Escribimos:

\[
f(x,y)
=
\frac{\sen y}{y}
\frac{x}{\sen x}.
\]

Como:

\[
\lim_{y\to0}\frac{\sen y}{y}=1
\]

y:

\[
\lim_{x\to0}\frac{x}{\sen x}=1,
\]

se obtiene:

\[
\boxed{
\lim_{(x,y)\to(0,0)}
\frac{x\sen y}{y\sen x}
=
1
}
\]

\section{Ejercicio 14}
\subsection{Solución}

Sea:

\[
f(x,y)
=
\frac{3x-12}{x^2-16y^2}.
\]

Según:

\[
y=1,
\]

se tiene:

\[
f(x,1)
=
\frac{3(x-4)}
{(x-4)(x+4)}
=
\frac3{x+4}.
\]

Por tanto:

\[
\lim_{x\to4}f(x,1)
=
\frac38.
\]

Según:

\[
x=4,
\]

se tiene:

\[
f(4,y)=0
\]

para \(y\neq1\), luego:

\[
\lim_{y\to1}f(4,y)=0.
\]

Como los resultados son distintos:

\[
\boxed{
\lim_{(x,y)\to(4,1)}
\frac{3x-12}{x^2-16y^2}
\text{ no existe}
}
\]

\section{Ejercicio 15}
\subsection{Solución}

Sea:

\[
f(x,y)
=
\sqrt{x}
\ln\left(
1+\sqrt{x^2-y^2}
+
\left|\frac yx\right|
\right).
\]

En su dominio, cerca del origen:

\[
x>0,
\qquad
|y|\leq x.
\]

Entonces:

\[
0\leq
\sqrt{x^2-y^2}
\leq x
\]

y:

\[
0\leq
\left|\frac yx\right|
\leq1.
\]

Por tanto:

\[
1
\leq
1+\sqrt{x^2-y^2}
+
\left|\frac yx\right|
\leq
2+x.
\]

Así:

\[
0
\leq
f(x,y)
\leq
\sqrt{x}\ln(2+x).
\]

Como:

\[
\lim_{x\to0^+}
\sqrt{x}\ln(2+x)
=
0,
\]

se concluye:

\[
\boxed{
\lim_{(x,y)\to(0,0)}
\sqrt{x}
\ln\left(
1+\sqrt{x^2-y^2}
+
\left|\frac yx\right|
\right)
=
0
}
\]

\section{Ejercicio 16}
\subsection{Solución}

Sea:

\[
f(x,y)
=
\frac{x^2y^2}
{x^2y^2+(x-y)^2}.
\]

Según:

\[
y=x,
\]

se obtiene:

\[
f(x,x)
=
\frac{x^4}{x^4}
=
1.
\]

Por tanto:

\[
\lim_{x\to0}f(x,x)=1.
\]

Según:

\[
y=0,
\]

se tiene:

\[
f(x,0)=0.
\]

Luego:

\[
\lim_{x\to0}f(x,0)=0.
\]

Como los valores son distintos:

\[
\boxed{
\lim_{(x,y)\to(0,0)}
\frac{x^2y^2}
{x^2y^2+(x-y)^2}
\text{ no existe}
}
\]

\section{Ejercicio 17}
\subsection{Solución}

Sea:

\[
f(x,y)
=
\left(
1-x\sen y,\,
3x^2e^{-y^2}
\right).
\]

Calculamos componente a componente:

\[
\lim_{(x,y)\to(0,0)}
(1-x\sen y)
=
1
\]

y:

\[
\lim_{(x,y)\to(0,0)}
3x^2e^{-y^2}
=
0.
\]

Por tanto:

\[
\boxed{
\lim_{(x,y)\to(0,0)}f(x,y)
=
(1,0)
}
\]

\section{Ejercicio 18}
\subsection{Solución}

Sea:

\[
t=x^2+y^2+z^2.
\]

Cuando:

\[
(x,y,z)\to(0,0,0),
\]

se tiene:

\[
t\to0.
\]

Por tanto:

\[
\frac{\sen(x^2+y^2+z^2)}
{x^2+y^2+z^2}
=
\frac{\sen t}{t}.
\]

Así:

\[
\boxed{
\lim_{(x,y,z)\to(0,0,0)}
\frac{\sen(x^2+y^2+z^2)}
{x^2+y^2+z^2}
=
1
}
\]

\section{Ejercicio 19}
\subsection{Solución}

Sea:

\[
f(x,y,z)
=
\frac{xy+yz}{x^2+y^2+z^2}.
\]

Según:

\[
(x,y,z)=(t,0,0),
\]

se obtiene:

\[
f(t,0,0)=0.
\]

Según:

\[
(x,y,z)=(t,t,0),
\]

se tiene:

\[
f(t,t,0)
=
\frac{t^2}{2t^2}
=
\frac12.
\]

Como se obtienen límites distintos:

\[
\boxed{
\lim_{(x,y,z)\to(0,0,0)}
\frac{xy+yz}{x^2+y^2+z^2}
\text{ no existe}
}
\]

\section{Ejercicio 20}
\subsection{Solución}

Sea:

\[
f(x,y)
=
\begin{cases}
\dfrac{\sen^2x\,\sen y}{x^2+y^2},
&
(x,y)\neq(0,0),\\[2mm]
0,
&
(x,y)=(0,0).
\end{cases}
\]

Fuera del origen es cociente de funciones continuas con denominador no nulo.

En el origen:

\[
|\sen x|\leq|x|,
\qquad
|\sen y|\leq|y|.
\]

Por tanto:

\[
\left|
\frac{\sen^2x\,\sen y}{x^2+y^2}
\right|
\leq
\frac{x^2|y|}{x^2+y^2}
\leq
|y|.
\]

Como:

\[
|y|\to0,
\]

se obtiene:

\[
\lim_{(x,y)\to(0,0)}f(x,y)=0=f(0,0).
\]

Por tanto:

\[
\boxed{
f\text{ es continua en todo }\R^2
}
\]

\section{Ejercicio 21}
\subsection{Solución}

Sea:

\[
f(x,y)
=
\begin{cases}
x\sen\left(\dfrac1{x^2+y^2}\right)
+
y\cos\left(\dfrac1{x^2+y^2}\right),
&
(x,y)\neq(0,0),\\[2mm]
0,
&
(x,y)=(0,0).
\end{cases}
\]

Para \((x,y)\neq(0,0)\):

\[
|f(x,y)|
\leq
|x|+|y|.
\]

Además:

\[
|x|+|y|
\leq
\sqrt2\sqrt{x^2+y^2}.
\]

Por tanto:

\[
\lim_{(x,y)\to(0,0)}f(x,y)=0.
\]

Como:

\[
f(0,0)=0,
\]

se concluye:

\[
\boxed{
f\text{ es continua en }(0,0)
}
\]

\section{Ejercicio 22}
\subsection{Solución}

Sea:

\[
f(x,y)
=
\begin{cases}
\dfrac{2xy^2+3x}{x^2+y^2},
&
(x,y)\neq(0,0),\\[2mm]
0,
&
(x,y)=(0,0).
\end{cases}
\]

Fuera del origen es continua.

En el origen, según:

\[
y=0,
\]

se obtiene:

\[
f(x,0)
=
\frac{3x}{x^2}
=
\frac3x.
\]

Este valor no posee límite cuando:

\[
x\to0.
\]

Por tanto:

\[
\boxed{
f\text{ no es continua en }(0,0)
}
\]

y:

\[
\boxed{
f\text{ es continua en }
\R^2\setminus\{(0,0)\}.
}
\]

\section{Ejercicio 23}
\subsection{Solución}

Sea:

\[
f(x,y)
=
\begin{cases}
\dfrac{x}{4x^2+y^2-1},
&
4x^2+y^2\neq1
\text{ y }(x,y)\neq(0,0),\\[2mm]
1,
&
4x^2+y^2=1
\text{ o }(x,y)=(0,0).
\end{cases}
\]

En los puntos donde:

\[
4x^2+y^2\neq1
\]

y:

\[
(x,y)\neq(0,0),
\]

la función es racional con denominador no nulo, luego es continua.

En el origen:

\[
\lim_{(x,y)\to(0,0)}
\frac{x}{4x^2+y^2-1}
=
0
\]

pero:

\[
f(0,0)=1.
\]

Por tanto, no es continua en el origen.

Si:

\[
4x_0^2+y_0^2=1
\]

y \(x_0\neq0\), el denominador tiende a \(0\) mientras el numerador tiende a \(x_0\neq0\), por lo que no existe un límite finito.

En los puntos:

\[
(0,1),
\qquad
(0,-1),
\]

tomando \(x=0\):

\[
f(0,y)=0
\]

para \(y\neq\pm1\), mientras que el valor asignado es \(1\).

Por tanto:

\[
\boxed{
f\text{ es continua exactamente en }
\R^2
\setminus
\left(
\{(0,0)\}
\cup
\{(x,y):4x^2+y^2=1\}
\right)
}
\]

\section{Ejercicio 24}
\subsection{Solución}

Sea:

\[
h(x,y,z)
=
\left(
3e^{x+y+z},
\sen(x^2+y^2),
1+x+y+z
\right).
\]

Sus tres componentes:

\[
h_1(x,y,z)=3e^{x+y+z},
\]

\[
h_2(x,y,z)=\sen(x^2+y^2),
\]

\[
h_3(x,y,z)=1+x+y+z
\]

son continuas en todo \(\R^3\).

Por tanto:

\[
\boxed{
h\text{ es continua en todo }\R^3
}
\]

\end{document}
